\section{A Reproducible LEO Case Study}
\label{sec:casestudy}

The case study comprises five experiments, summarized in Table~\ref{tab:expsummary};
each isolates one question, from the bare mechanism to scale and hardware noise, and each
is reproducible from seeded runs in the companion code.

\begin{table}[t]
\centering
\caption{The five case-study experiments and what each isolates.}
\label{tab:expsummary}
\renewcommand{\arraystretch}{1.25}
\footnotesize
\begin{tabular}{@{}p{0.07\columnwidth}p{0.40\columnwidth}p{0.40\columnwidth}@{}}
\toprule
Exp & Question & Finding \\
\midrule
A & Does the quantum walk spread faster than diffusion? & Yes, exactly: ballistic
$\Var\!\sim\!t^2$ vs diffusive $\Var\!\sim\!t$ \\
B & Operator comparison on the LEO task & global $\gg$ local ($\sim\!\globalgain\times$);
QW ties heat within seed noise \\
B2 & Does any quantum-walk variant win? & adjacency walk is best, edge on far nodes,
within noise on MAE \\
C & Two evaluation pitfalls & scale-coupled norm and single-seed both fabricate
advantages \\
D & Does the mechanism survive device noise? & no: depolarizing noise decoheres the
walk toward uniform \\
E & Does the ranking hold at scale? & quantum-inspired walk overtakes heat at $1584$;
the ballistic edge grows with diameter \\
\bottomrule
\end{tabular}
\end{table}

\subsection{Setup}
We instantiate a Walker-delta LEO constellation with pure-NumPy circular-orbit
propagation, so the entire topology is reproducible from parameters alone. Inter
satellite links follow the standard $+$Grid pattern with a polar cutoff; restricting
to the largest connected component yields a snapshot graph. The case-study scale is a
$264$-satellite shell; a smaller $66$-satellite shell is used where noted for the
seed study.

\textbf{Task.} For a randomly chosen destination $v_d$, each node receives the input
feature $[\mathbf{1}\{i=v_d\},\,\deg_i/\max\deg]$, a single seeded signal, and the
target is the graph-geodesic hop distance $y_i=d_G(i,v_d)$ normalized by the source
eccentricity. The model must propagate the seed across the graph to recover the
distance field (Fig.~\ref{fig:leotask}). This task is deliberately decoupled from any
particular routing or traffic-engineering system: it isolates the propagation question.

\textbf{Reproducibility details.} The constellations are Walker-delta shells: a
$66$-satellite shell ($6$ planes of $11$, inclination $86.4^\circ$, altitude $780$~km), a
$264$-satellite shell ($24\times 11$, $53^\circ$, $550$~km), and the Starlink shell-1
geometry ($1584$ satellites, $72\times 22$, $53^\circ$, $550$~km) for the scale test.
Positions follow pure circular-orbit propagation in an Earth-centered inertial frame, and
inter-satellite links use the $+$Grid pattern with a $70^\circ$ polar cutoff and the
counter-rotating seam cut. Each instance draws a random snapshot time and a random
destination; the Laplacian is eigendecomposed once and shared across operators. Every
model is a three-to-four-layer network of width $32$ trained with Adam for a few hundred
epochs; all randomness (initialization, snapshot sampling, destination choice) is seeded,
and the seeds, hyperparameters, and figure scripts are in the companion repository so each
number here is regenerable. We deliberately avoid any tuning that differs between
operators, so the comparison stays param-matched.

\textbf{Protocol (a rigor battery).} We follow the evaluation discipline that
robustness and dependability venues now expect, where the trustworthiness of the
\emph{evaluation} is treated as seriously as the headline number. Four rules are fixed
before any result is read. (i) \emph{Param-matched operators}: all models share depth,
width, and weight shapes (Section~\ref{sec:buildit}), so a gap is attributable to the
operator and not to capacity. (ii) \emph{Multiple seeds}: every reported number is a
mean and standard deviation over at least three seeds (five for the operator
comparison, eight for the seed study), never a single run; a ranking that lies inside
the combined spread is reported as a tie, not a result. (iii) \emph{Scale-free summary
metric}: alongside mean absolute error (MAE) we report a far-node
area-under-the-error-curve (AURC), the error integrated over a difficulty axis
$\rho\in[0,1]$ where bin $\rho$ keeps nodes with $d_G(i,v_d)\ge\rho\,\mathrm{ecc}(v_d)$.
The AURC is a single number, robust to picking a flattering far-node threshold and to
the graph scale. (iv) \emph{Design sweep before a verdict}: we vary the obvious
quantum-walk design choices (Exp~B2) so that any conclusion is about the mechanism, not
one implementation. Every number in the tables and figures is regenerated from seeded
runs by the companion code, so the paper and the released artifact cannot drift apart.

\begin{figure}[t]
\centering
\resizebox{0.85\columnwidth}{!}{%
\begin{tikzpicture}[node distance=5mm]
% a small grid-like LEO graph
\foreach \r in {0,1,2} {
  \foreach \c in {0,1,2,3} {
    \node[circle, draw=steel, fill=white, inner sep=1.4pt] (n\r\c) at (\c*1.1, \r*1.1) {};
  }
}
% intra/inter links
\foreach \r in {0,1,2} \foreach \c in {0,1,2} {
  \pgfmathtruncatemacro\cn{\c+1}
  \draw[steel] (n\r\c) -- (n\r\cn);
}
\foreach \r in {0,1} \foreach \c in {0,1,2,3} {
  \pgfmathtruncatemacro\rn{\r+1}
  \draw[steel] (n\r\c) -- (n\rn\c);
}
% destination seed
\node[circle, draw=amber!80!black, fill=palegold, inner sep=1.8pt] at (n00) {};
\node[font=\scriptsize, text=ink, below=1mm of n00] {$v_d$ (seed $=1$)};
% target field shading hint
\node[font=\scriptsize, text=ink, above=1mm of n23] {far node: $y=d_G(\cdot,v_d)$};
\end{tikzpicture}}
\caption{The node-potential task. A single seeded signal at the destination $v_d$ must
be propagated across the dynamic LEO graph to recover the geodesic-distance field. The
task rewards long-range operators and exposes a local GCN's $K$-hop horizon.}
\label{fig:leotask}
\end{figure}


\subsection{Mechanism: ballistic vs diffusive (Exp A)}
We first verify the mechanism in isolation. Releasing a unit impulse at the center of
a path graph and propagating it, we measure the spatial second moment versus
propagation time for both operators. Fig.~\ref{fig:expA} confirms the theory exactly:
the heat kernel spreads diffusively ($\Var\sim t$) and the quantum walk ballistically
($\Var\sim t^2$), matching \eqref{eq:diffusive} and \eqref{eq:ballistic}. A log-log
fit recovers slopes of $1.00$ and $2.00$. This is the clean, theory-backed sense in
which the quantum operator reaches farther per unit time.

\begin{figure}[t]
\centering
\resizebox{\columnwidth}{!}{%
\begin{tikzpicture}
\begin{axis}[
  width=8.4cm, height=5.7cm,
  xlabel={propagation time $t$}, ylabel={spread variance},
  xmin=0, xmax=8, ymin=0, ymax=135,
  legend pos=north west, legend cell align=left,
  grid=major, grid style={gray!18},
  tick label style={font=\footnotesize}, label style={font=\small},
  legend style={font=\small, fill=white, fill opacity=0.85, draw=gray!40},
  axis line style={gray!50},
]
% quantum walk: ballistic, area fill (baseline corners close the region)
\addplot[draw=amber!80!black, very thick, mark=square*, mark size=1.4pt,
         fill=palegold, fill opacity=0.45] coordinates {
  (0,0)(1,2)(2,8)(3,18)(4,32)(5,50)(6,72)(7,98)(8,128)(8,0)} \closedcycle;
\addlegendentry{quantum walk $|e^{-it\Lap}|^2\ (\sim t^2)$}
% heat: diffusive, drawn on top
\addplot[draw=steel, very thick, mark=*, mark size=1.4pt,
         fill=paleblue, fill opacity=0.6] coordinates {
  (0,0)(1,2)(2,4)(3,6)(4,8)(5,10)(6,12)(7,14)(8,16)(8,0)} \closedcycle;
\addlegendentry{heat $e^{-t\Lap}\ (\sim t)$}
\node[font=\footnotesize\bfseries, text=amber!75!black, rotate=58]
  at (axis cs:6.0,86) {ballistic};
\node[font=\footnotesize\bfseries, text=steel]
  at (axis cs:6.7,21) {diffusive};
\end{axis}
\end{tikzpicture}}
\caption{Exp~A, mechanism. Spatial second moment of a propagated impulse on a path
graph versus propagation time. Heat diffusion grows linearly ($\Var\sim t$) and the
quantum walk ballistically ($\Var\sim t^2$); a log-log fit gives slopes $1.00$ and
$2.00$, matching \eqref{eq:diffusive} and \eqref{eq:ballistic}. The shaded gap is the
quadratic-versus-linear reach advantage.}
\label{fig:expA}
\end{figure}


\subsection{Operator comparison (Exp B)}
We then train the three operators on the LEO task over five seeds. The results are
summarized in Table~\ref{tab:operators} and Fig.~\ref{fig:expB}. Two findings stand
out, and we state both honestly.

\textit{Global beats local, by a lot.} Both global operators reduce error roughly
$\globalgain\times$ relative to the local GCN (MAE $\bMAEheat$ and $\bMAEqw$ versus
$\bMAEgcn$; the gap is even larger on the far-node AURC, $\bAURCheat$ and $\bAURCqw$
versus $\bAURCgcn$). This is the robust, unambiguous result of the study: a $K$-hop
GCN cannot cover a graph whose diameter exceeds $K$, while a global operator can.

\textit{Quantum-inspired ties classical, within noise.} The quantum-inspired walk and
the classical heat kernel are statistically indistinguishable on overall MAE
($\bMAEqw\pm\bSTDqw$ versus $\bMAEheat\pm\bSTDheat$), and the quantum-inspired operator
is the noisier of the two. The one place it shows a consistent edge is the far-node
AURC ($\bAURCqw$ versus $\bAURCheat$), which is exactly where ballistic spreading
should help: on the hardest, most distant nodes. The edge is small and stays within
the seed-to-seed spread, so \emph{at this $264$-satellite scale} we do not claim a
quantum advantage; we report a theory-consistent tendency that a single seed would have
over- or under-stated, and that the scale experiment (Section~\ref{sec:expe}) will turn
into a clear win at larger diameter.

\begin{table}[t]
\centering
\caption{Operator comparison on the LEO node-potential task ($264$ satellites,
$5$ seeds, param-matched). Lower is better.}
\label{tab:operators}
\renewcommand{\arraystretch}{1.2}
\begin{tabular}{@{}lccc@{}}
\toprule
Operator & MAE & far-node AURC & params \\
\midrule
GCN (local)            & $\bMAEgcn\pm\bSTDgcn$  & $\bAURCgcn$ & $3301$ \\
Heat (global)          & $\bMAEheat\pm\bSTDheat$ & $\bAURCheat$ & $3301$ \\
Quantum-inspired (global) & $\bMAEqw\pm\bSTDqw$  & $\bAURCqw$  & $3301$ \\
\bottomrule
\end{tabular}
\end{table}

\begin{figure}[t]
\centering
\resizebox{\columnwidth}{!}{%
\begin{tikzpicture}
\begin{axis}[
  width=8.4cm, height=5.7cm,
  ybar, bar width=12pt,
  symbolic x coords={GCN, Heat, QI},
  xtick=data,
  ymin=0, ymax=0.24,
  ylabel={error (lower is better)},
  legend pos=north east, legend cell align=left,
  ymajorgrids, grid style={gray!18},
  tick label style={font=\footnotesize}, label style={font=\small},
  legend style={font=\footnotesize, fill=white, fill opacity=0.85, draw=gray!40},
  axis line style={gray!50},
  error bars/y dir=both, error bars/y explicit,
  enlarge x limits=0.28,
]
\addplot[draw=steel!80!black, fill=steel!35] coordinates {
  (GCN,0.163) +- (0,0.016)
  (Heat,0.069) +- (0,0.005)
  (QI,0.071) +- (0,0.009)};
\addlegendentry{MAE}
\addplot[draw=amber!80!black, fill=palegold] coordinates {
  (GCN,0.202) (Heat,0.085) (QI,0.077)};
\addlegendentry{far-node AURC}
\end{axis}
\end{tikzpicture}}
\caption{Exp~B, operator comparison on the LEO node-potential task ($264$ satellites,
$5$ seeds, param-matched). Global operators (Heat, QI) beat the local GCN by roughly
$\globalgain\times$. Heat and the quantum-inspired (QI) walk tie on MAE within the seed
spread (error bars); QI shows a small, theory-consistent edge only on the far-node AURC.}
\label{fig:expB}
\end{figure}


\textit{Why the seeds matter.} Table~\ref{tab:perseed} lists the five per-seed runs
behind the aggregate. They make the within-noise verdict concrete: the quantum-inspired
walk wins on seeds $0$ and $3$ and loses on seeds $1$, $2$, and $4$ (seed $1$ a near
tie), while the local GCN loses on every seed by a wide margin. Reporting seed $0$ alone
would have advertised a clean quantum win ($0.068$ versus $0.077$); reporting seed $2$
alone would have shown a clean quantum loss ($0.082$ versus $0.068$). Neither is the
truth. Only the
distribution, $\bMAEqw\pm\bSTDqw$ versus $\bMAEheat\pm\bSTDheat$, supports an honest
statement, and that statement is ``tie.'' This single table is the strongest argument in
the tutorial for never trusting a single seed.

\begin{table}[t]
\centering
\caption{Per-seed MAE behind Table~\ref{tab:operators} (LEO task, $264$ satellites).
Bold marks the better (lower) per-seed value. The quantum-inspired walk and heat trade
wins across seeds; the local GCN loses every seed. The ranking inside the seed spread is
a tie, not a result.}
\label{tab:perseed}
\renewcommand{\arraystretch}{1.15}
\begin{tabular}{@{}lccccc|c@{}}
\toprule
Operator & s0 & s1 & s2 & s3 & s4 & mean \\
\midrule
GCN  & .148 & .176 & .168 & .181 & .140 & $\bMAEgcn$ \\
Heat & .077 & \textbf{.063} & \textbf{.068} & .063 & \textbf{.072} & $\bMAEheat$ \\
QI   & \textbf{.068} & .063 & .082 & \textbf{.062} & .080 & $\bMAEqw$ \\
\bottomrule
\end{tabular}
\end{table}

\subsection{Trying to rescue the advantage (Exp B2)}
Before concluding, we probe whether a different quantum-walk design recovers a clear
win: a walk on the adjacency $e^{-it\Adj}$, a walk on the normalized Laplacian
$e^{-it\Lap^{\mathrm{sym}}}$, and a multi-time readout that concatenates several
learnable propagation times. The first two are param-matched; the multi-time variant
is given extra capacity, so it is a generous upper bound for the family.
Table~\ref{tab:variants} reports the outcome, and the sweep matters.
\bbVariantNote

The methodological lesson is the point. Had we stopped at the Laplacian walk of
Exp~B, we would have concluded that the quantum-inspired operator merely ties the
classical one. The sweep shows that conclusion was about one implementation, not the
mechanism: a different and equally natural choice, the adjacency walk, is in fact the
best operator we tested, at matched parameters and with the lowest variance. At this
$264$-satellite scale the overall-MAE margin over heat
($\bbAdjMAE\pm\bbAdjMAEstd$ versus $\bbHeatMAE$) still stays within the seed spread; the
honest verdict moves from ``ties'' to ``the best variant, with a consistent edge on the
hardest nodes,'' and only a design sweep could have told the difference. The scale
experiment (Section~\ref{sec:expe}) then shows this edge widening into a clear win at
$1584$ satellites.

\begin{table}[t]
\centering
\caption{Quantum-walk variant probe on the LEO task ($264$ satellites, $5$ seeds).
Lower is better. The multi-time variant is capacity-favoured.}
\label{tab:variants}
\renewcommand{\arraystretch}{1.2}
\begin{tabular}{@{}lccc@{}}
\toprule
Variant & MAE & far-node AURC & params \\
\midrule
Heat (reference)                 & $\bbHeatMAE$ & $\bbHeatAURC$ & $3301$ \\
QW on $\Lap$ (baseline)          & $\bbLapMAE$  & $\bbLapAURC$  & $3301$ \\
\textbf{QW on $\Adj$}            & $\mathbf{\bbAdjMAE}$ & $\mathbf{\bbAdjAURC}$ & $3301$ \\
QW on $\Lap^{\mathrm{sym}}$      & $\bbNormMAE$ & $\bbNormAURC$ & $3301$ \\
QW multi-time (higher capacity)  & $\bbMultiMAE$ & $\bbMultiAURC$ & $9581$ \\
\bottomrule
\end{tabular}
\end{table}

\begin{figure}[t]
\centering
\resizebox{\columnwidth}{!}{%
\begin{tikzpicture}
\begin{axis}[
  width=8.6cm, height=5.7cm,
  ybar, bar width=7pt,
  symbolic x coords={heat, qw-L, qw-A, qw-Lsym, qw-multi},
  xtick=data, x tick label style={font=\scriptsize, rotate=18, anchor=east},
  ymin=0.05, ymax=0.099,
  ylabel={error (lower is better)},
  legend cell align=left,
  legend style={font=\scriptsize, at={(0.02,1.02)}, anchor=south west,
                legend columns=2, fill=white, draw=gray!40},
  ymajorgrids, grid style={gray!18},
  tick label style={font=\footnotesize}, label style={font=\small},
  axis line style={gray!50},
  enlarge x limits=0.13,
]
\addplot[draw=steel!80!black, fill=steel!35] coordinates {
  (heat,0.069) (qw-L,0.071) (qw-A,0.066) (qw-Lsym,0.071) (qw-multi,0.069)};
\addlegendentry{MAE}
\addplot[draw=amber!80!black, fill=palegold] coordinates {
  (heat,0.085) (qw-L,0.077) (qw-A,0.075) (qw-Lsym,0.085) (qw-multi,0.082)};
\addlegendentry{far-node AURC}
% highlight the best variant
\node[star, star points=5, star point ratio=2.3, fill=amber, draw=amber!70!black,
      inner sep=1.3pt] at (axis cs:qw-A,0.0905) {};
\node[font=\scriptsize\bfseries, text=amber!75!black] at (axis cs:qw-A,0.087) {best};
\end{axis}
\end{tikzpicture}}
\caption{Exp~B2, quantum-walk variant probe ($264$ satellites, $5$ seeds). The base
matrix of the walk matters: the adjacency walk (qw-A, starred) is the best operator on
both metrics at matched parameters, edging the classical heat kernel on the far-node
AURC, while the Laplacian (qw-L) and normalized-Laplacian (qw-Lsym) walks tie heat and
the higher-capacity multi-time variant does not improve on qw-A.}
\label{fig:expB2}
\end{figure}


\subsection{From quantum-inspired to true quantum: NISQ noise (Exp D)}
The case study uses the quantum-walk kernel as a \emph{classical} layer. To see what is
lost on real hardware, we run the continuous-time quantum walk of Exp~A as an actual
quantum circuit (a PennyLane mixed-state simulator on an $8$-node path) and inject
per-qubit depolarizing noise of strength $p$. Fig.~\ref{fig:expD} shows the coherent,
ballistic structure decohering toward the uniform mixed state as $p$ grows: the $L_1$
distance to uniform falls from $\dLoneClean$ to $\dLoneNoisy$, and the spread variance
drops from its clean value $\dVarClean$ toward the featureless uniform level $5.25$. The
ballistic mechanism that Exp~A verified exactly is therefore fragile to noise, which is
the concrete reason the runnable path today is the classical quantum-inspired surrogate,
not a NISQ implementation. A fuller hardware study (noise-aware training, larger graphs)
is left to future work (Section~\ref{sec:roadmap}).

\begin{figure}[t]
\centering
\resizebox{\columnwidth}{!}{%
\begin{tikzpicture}
\begin{axis}[
  width=8.4cm, height=5.7cm,
  xlabel={depolarizing noise $p$}, ylabel={value},
  xmin=0, xmax=0.5, ymin=0, ymax=1.05,
  legend pos=south west, legend cell align=left,
  grid=major, grid style={gray!18},
  tick label style={font=\footnotesize}, label style={font=\small},
  legend style={font=\scriptsize, fill=white, fill opacity=0.9, draw=gray!40},
  axis line style={gray!50},
]
% L1 distance to the uniform (decohered) state: coherence collapses with noise
\addplot[amber!80!black, very thick, mark=square*, mark size=1.4pt] coordinates {
  (0,0.569)(0.02,0.525)(0.05,0.463)(0.1,0.370)(0.2,0.242)(0.3,0.160)(0.5,0.067)};
\addlegendentry{$L_1$ to uniform (coherence)}
% spread variance normalized to its clean value, for shape comparison
\addplot[steel, very thick, mark=*, mark size=1.4pt] coordinates {
  (0,1.000)(0.02,0.986)(0.05,0.967)(0.1,0.936)(0.2,0.881)(0.3,0.835)(0.5,0.770)};
\addlegendentry{spread var (normalized)}
\draw[gray!60, dashed] (axis cs:0,0.741) -- (axis cs:0.5,0.741);
\node[font=\scriptsize, text=gray!55, anchor=south west] at (axis cs:0.10,0.748)
  {uniform-state variance};
\end{axis}
\end{tikzpicture}}
\caption{Exp~D, NISQ noise erodes the ballistic mechanism. A continuous-time quantum
walk on an $8$-node path is run as a PennyLane mixed-state circuit; per-qubit
depolarizing noise of strength $p$ drives the coherent walk toward the uniform mixed
state. The $L_1$ distance to uniform falls from $\dLoneClean$ at $p{=}0$ to $\dLoneNoisy$
at $p{=}0.5$, and the spread variance decays from its clean value ($\dVarClean$) toward
the featureless uniform-state level ($5.25$, dashed). This is the near-term caveat behind
preferring the classical quantum-inspired surrogate.}
\label{fig:expD}
\end{figure}


\subsection{The quantum-inspired edge emerges at scale (Exp E)}
\label{sec:expe}
To test the operator ranking beyond the $264$-satellite scale, we repeat the comparison,
at the same training budget, on a Starlink shell-1 graph of $1584$ satellites, the largest
at which a full eigendecomposition of $\Lap$ remains affordable on a CPU.
Table~\ref{tab:scale} and Fig.~\ref{fig:scale} report MAE across the three shells, and the
trend is the most interesting result in the study. Global operators keep their large lead
over the local GCN. More tellingly, the quantum-inspired and classical global operators
\emph{change places as the graph grows}: at $66$ satellites heat wins
($\cTwoHeat$ vs $\cTwoQw$), at $264$ they tie ($\bMAEheat$ vs $\bMAEqw$), and at $1584$
the quantum-inspired walk is clearly best, $\eMAEqw\!\pm\!\eMAEqwStd$ versus heat's
$\eMAEheat\!\pm\!\eMAEheatStd$, winning two of three seeds with markedly lower variance,
and far ahead on the far-node metric ($\eAURCqw$ versus $\eAURCheat$). This is exactly
what the ballistic mechanism predicts: an operator whose reach grows as $t$ rather than
$\sqrt{t}$ should help most where geodesics are longest, and diameter grows with the
constellation. We temper the claim only by the three-seed budget at this scale; within
that limit the trend is monotone and theory-consistent, and it directly answers the
external-validity concern of Section~\ref{sec:threats}.

\begin{table}[t]
\centering
\caption{Operator MAE across constellation scale (lower is better), same training budget
at every scale. The quantum-inspired walk goes from losing ($66$) to tying ($264$) to
winning ($1584$) against the classical heat kernel, as the ballistic mechanism predicts.}
\label{tab:scale}
\renewcommand{\arraystretch}{1.2}
\begin{tabular}{@{}lccc@{}}
\toprule
Operator & $66$ sats & $264$ sats & $1584$ sats \\
\midrule
GCN (local)               & $\cTwoGcn$  & $\bMAEgcn$  & $\eMAEgcn$ \\
Heat (global)             & $\cTwoHeat$ & \textbf{$\bMAEheat$} & $\eMAEheat$ \\
Quantum-inspired (global) & $\cTwoQw$   & $\bMAEqw$   & \textbf{$\eMAEqw$} \\
\bottomrule
\end{tabular}
\end{table}

\begin{figure}[t]
\centering
\resizebox{\columnwidth}{!}{%
\begin{tikzpicture}
\begin{axis}[
  width=8.4cm, height=5.7cm,
  xlabel={constellation size (satellites)}, ylabel={MAE (lower is better)},
  symbolic x coords={66, 264, 1584}, xtick=data,
  ymin=0.04, ymax=0.23,
  legend pos=north west, legend cell align=left,
  grid=major, grid style={gray!18},
  tick label style={font=\footnotesize}, label style={font=\small},
  legend style={font=\scriptsize, fill=white, fill opacity=0.9, draw=gray!40},
  axis line style={gray!50},
]
\addplot[gray!55, very thick, mark=triangle*, mark size=2pt] coordinates {
  (66,0.102) (264,0.163) (1584,0.208)};
\addlegendentry{GCN (local)}
\addplot[steel, very thick, mark=*, mark size=2pt] coordinates {
  (66,0.064) (264,0.069) (1584,0.168)};
\addlegendentry{Heat (classical global)}
\addplot[amber!80!black, very thick, mark=square*, mark size=2pt] coordinates {
  (66,0.085) (264,0.071) (1584,0.148)};
\addlegendentry{Quantum-inspired walk}
\node[font=\scriptsize\bfseries, text=amber!75!black, align=center]
  at (axis cs:1584,0.118) {QI best};
\end{axis}
\end{tikzpicture}}
\caption{Exp~E, the quantum-inspired edge emerges at scale. Operator MAE versus
constellation size at a fixed training budget. The quantum-inspired walk (amber) starts
worse than the classical heat kernel (blue) at $66$ satellites, ties at $264$, and
overtakes it at $1584$, where geodesics are longest and the ballistic ($\Var\!\sim\!t^2$)
reach of the walk pays off; the local GCN trails at every scale.}
\label{fig:scale}
\end{figure}


\subsection{Discussion: reading the findings together}
The case study tells a consistent story across its parts, and the consistency is what
makes it trustworthy. The mechanism experiment (Exp~A) confirms, exactly and without
training, that the quantum walk reaches farther per unit time than diffusion. The
operator comparison (Exp~B) shows that reach is decisive against a local operator: the
GCN, capped at its hop horizon, loses by roughly $\globalgain\times$, while at the
$264$-satellite scale heat and the Laplacian walk are statistically tied. The variant
probe (Exp~B2) shows the tie is implementation-dependent, the adjacency walk edging heat
on the far-node metric. The scale experiment (Exp~E) then supplies the unifying thread:
the quantum-inspired walk loses at $66$ satellites, ties at $264$, and \emph{wins} at
$1584$, exactly the monotone, diameter-driven trend the ballistic mechanism predicts. The
four views do not merely fail to contradict one another; they line up along a single
axis, the longer the geodesics, the more the $\Var\!\sim\!t^2$ reach of the walk matters.
That alignment is the signature of a real mechanism rather than a lucky configuration.

The same consistency bounds how far we may push the claim. The scale trend rests on a
three-seed budget at $1584$ satellites, and at small and medium scale the quantum-inspired
operator is at best a tie and is noisier than heat. The honest synthesis for a
practitioner is therefore: a global operator is mandatory; the classical heat kernel is an
excellent default at small-to-medium diameter; and the quantum-inspired walk is worth
adopting where the graph is large and long-range, the regime where our largest-scale
experiment shows it is the best operator. That is a more useful and more durable message
than a blanket advantage would have been, and it is the message the data actually
supports.
